\documentclass[a4paper, serif, 10pt]{moderncv}

%% moderncv
% color and style
\moderncvcolor{black}
\moderncvstyle{banking}

%% page layout/margins
\usepackage[top=2.5cm, bottom=2.5cm, left=1.5cm, right=1.5cm]{geometry}

\nopagenumbers{}

%% Babel config for Indonesian (Bahasa Indonesia) language
% ref:
% - https://latex3.github.io/babel/guides/locale-indonesian.html
\usepackage[indonesian]{babel}

%% fontspec
\usepackage{fontspec}

\IfFontExistsTF{Linux Libertine}{
  \setmainfont[Ligatures={TeX, Common}, Numbers=OldStyle]{Linux Libertine}
}{}
\IfFontExistsTF{Linux Libertine O}{
  \setmainfont[Ligatures={TeX, Common}, Numbers=OldStyle]{Linux Libertine O}
}{}

%% CTeX
% CJK support, used to show CJK characters in the resume
%
% - heading=false: disable ctex heading
% - punct=kaiming: use kaiming punctuations styles for CJK
% - scheme=plain: use plain scheme, do not override \normalsize font size
% - space=auto: space settings for CJK characters
%
% CJK fonts are resolved per locale language, e.g. Noto Serif/Sans CJK JP
% for Japanese resumes, with fallback to the next region if not installed
%
% ref:
% - http://ctan.mirrorcatalogs.com/language/chinese/ctex/ctex.pdf
\usepackage[UTF8, heading=false, punct=kaiming, scheme=plain, space=auto]{ctex}

\IfFontExistsTF{Noto Serif CJK SC}{
  \setCJKmainfont{Noto Serif CJK SC}
}{}
\IfFontExistsTF{Noto Sans CJK SC}{
  \setCJKsansfont{Noto Sans CJK SC}
}{}

%% line spacing
\usepackage{setspace}
\setstretch{1.125}

%% URL styling - use normal text instead of monospace
\urlstyle{same}

%% Patch \httplink and \httpslink to handle full URLs
% The original moderncv macros always prepend http:// or https:// to URLs.
% This causes issues when the URL already has a protocol (e.g., https://example.com)
% resulting in malformed URLs like https://https://example.com.
% This patch checks if the URL already contains "://" and uses it directly if so.
% Note: We use \str_set:Nx (x-expansion) to expand macros like \@homepage first.
\ExplSyntaxOn
\str_new:N \g_yamlresume_url_str

\RenewDocumentCommand{\httpslink}{O{}m}{
  \str_set:Nx \g_yamlresume_url_str {#2}
  \str_if_in:NnTF \g_yamlresume_url_str {://}
    { \url{#2} }
    { \url{https://#2} }
}

\RenewDocumentCommand{\httplink}{O{}m}{
  \str_set:Nx \g_yamlresume_url_str {#2}
  \str_if_in:NnTF \g_yamlresume_url_str {://}
    { \url{#2} }
    { \url{http://#2} }
}
\ExplSyntaxOff

\name{Rizky Pratama}{}
\title{QA Engineer dengan pengalaman dalam pengujian manual dan otomasi}
\phone[mobile]{+62 812-3456-7890}
\email{rizky.pratama.qa@gmail.com}
\homepage{https://www.rizkypratama.dev}

\address{Jl. Kebon Jeruk Raya No. 45, Jakarta Barat, DKI Jakarta, Indonesia, 11530}{}{}

\extrainfo{{\small \faLinkedin}\ \href{https://www.linkedin.com/in/rizkypratama}{@rizkypratama} {} {} {} • {} {} {} 
{\small \faGithub}\ \href{https://github.com/rizkypratama}{@rizkypratama}}

\begin{document}

\maketitle

\section{Informasi Dasar}

\cvline{}{QA Engineer dengan pengalaman lebih dari 7 tahun dalam memastikan kualitas perangkat lunak melalui pengujian manual dan otomasi. Terbiasa merancang strategi pengujian, menulis test case, melakukan regression testing, serta mengelola defect menggunakan JIRA.
%
\begin{itemize}
\item Berpengalaman mengotomasi pengujian web dan API menggunakan Selenium, Cypress, Playwright, dan Postman.
\item Mampu berkolaborasi dengan tim product, developer, dan DevOps dalam lingkungan Agile/Scrum.
\item Memiliki pemahaman kuat tentang SDLC, STLC, risk-based testing, dan continuous integration.
\end{itemize}}
\section{Pendidikan}

\cventry{Sep 2014–Jul 2018}
        {Sarjana, Teknik Informatika, Nilai: 3.65}
        {Universitas Bina Nusantara}
        {\url{https://www.binus.ac.id}}
        {}
        {\begin{itemize}
\item Mempelajari rekayasa perangkat lunak, basis data, dan pengujian perangkat lunak secara mendalam.
\item Menyelesaikan proyek akhir tentang otomasi pengujian aplikasi berbasis web.
\item Aktif dalam organisasi kemahasiswaan dan kompetisi teknologi.
\end{itemize}
\textbf{Mata Kuliah}: Rekayasa Perangkat Lunak, Basis Data, Pemrograman Web, Jaringan Komputer, Sistem Operasi, Pengujian Perangkat Lunak}
\section{Pengalaman Kerja}

\cventry{Mar 2021–Sekarang}
        {Senior QA Engineer}
        {PT Solusi Digital Nusantara}
        {\url{https://www.solusidigital.co.id}}
        {}
        {\begin{itemize}
\item Memimpin tim QA beranggotakan 4 orang dalam pengujian platform SaaS dengan lebih dari 50.000 pengguna aktif.
\item Merancang strategi pengujian end-to-end, test plan, dan test case untuk fitur baru serta regression suite.
\item Mengotomasi 300+ skenario pengujian menggunakan Cypress dan Playwright, menurunkan waktu regression testing sebesar 60\%.
\item Mengintegrasikan pengujian otomasi ke pipeline CI/CD Jenkins dan GitLab CI.
\item Berkolaborasi dengan product manager dan developer untuk mempercepat siklus rilis dua mingguan.
\item Melakukan mentoring dan knowledge sharing untuk QA junior.
\end{itemize}
\textbf{Kata Kunci}: Otomasi Pengujian, Cypress, Playwright, CI/CD, Kepemimpinan}

\cventry{Agu 2018–Feb 2021}
        {QA Engineer}
        {PT Karya Teknologi Indonesia}
        {\url{https://www.karyateknologi.co.id}}
        {}
        {\begin{itemize}
\item Melakukan pengujian manual dan otomasi pada aplikasi web dan mobile.
\item Menyusun test case berdasarkan user story dan acceptance criteria.
\item Melakukan API testing menggunakan Postman dan validasi data menggunakan SQL.
\item Melaporkan dan melacak defect di JIRA hingga tervalidasi.
\item Berkontribusi pada peningkatan proses QA dan dokumentasi test.
\end{itemize}
\textbf{Kata Kunci}: Pengujian Manual, Selenium, Postman, SQL, Agile}
\section{Bahasa}

\cvline{Indonesia}{Bahasa Ibu / Bilingual \hfill \textbf{Kata Kunci}: Bahasa Ibu}
\cvline{Inggris}{Tingkat Profesional Penuh \hfill \textbf{Kata Kunci}: TOEFL 580}
\section{Keahlian}

\cvline{Pengujian Manual}{Ahli \hfill \textbf{Kata Kunci}: Perancangan Test Case, Regression Testing, Exploratory Testing, User Acceptance Testing}
\cvline{Otomasi Pengujian}{Mahir \hfill \textbf{Kata Kunci}: Selenium, Cypress, Playwright, Appium, Katalon Studio}
\cvline{API \& Database}{Mahir \hfill \textbf{Kata Kunci}: Postman, REST API, SQL, MongoDB, GraphQL}
\cvline{Alat \& Proses}{Mahir \hfill \textbf{Kata Kunci}: JIRA, TestRail, Git, Jenkins, Agile/Scrum}
\section{Penghargaan}

\cventry{Des 2022}
        {PT Solusi Digital Nusantara}
        {QA Team of the Year}
        {}
        {}
        {Diberikan atas kontribusi dalam meningkatkan stabilitas rilis dan menurunkan defect produksi sebesar 35\% melalui inisiatif otomasi regression testing.}
\section{Sertifikat}

\cventry{Jun 2019}
        {ISTQB}
        {ISTQB Certified Tester Foundation Level}
        {\url{https://www.istqb.org/certifications/certified-tester-foundation-level}}
        {}
        {}

\cventry{Mar 2022}
        {Scrum Alliance}
        {Certified ScrumMaster}
        {\url{https://www.scrumalliance.org/get-certified/scrum-master-track/certified-scrummaster}}
        {}
        {}

\cventry{Sep 2021}
        {Cypress.io}
        {Cypress End-to-End Testing}
        {\url{https://docs.cypress.io/guides/end-to-end-testing/writing-your-first-end-to-end-test}}
        {}
        {}
\section{Publikasi}

\cventry{Agu 2023}
        {Meningkatkan Efektivitas Regression Testing melalui Otomasi Berbasis Risiko}
        {Jurnal Teknologi Informasi Indonesia}
        {\url{https://www.jti.ac.id/index.php/jti/article/view/12345}}
        {}
        {\begin{itemize}
\item Membahas pendekatan otomasi regression testing berbasis risiko untuk aplikasi web berskala besar.
\item Menunjukkan pengurangan waktu pengujian dan peningkatan cakupan test case.
\end{itemize}}
\section{Proyek}

\cventry{Jan 2022–Nov 2022}
        {Proyek otomasi pengujian end-to-end untuk platform e-commerce menggunakan Cypress dan TypeScript.}
        {Otomasi Pengujian Platform E-Commerce}
        {\url{https://github.com/rizkypratama/ecommerce-qa-automation}}
        {}
        {\begin{itemize}
\item Mengotomasi alur kritis seperti login, pencarian produk, keranjang belanja, dan checkout.
\item Menerapkan Page Object Model dan data-driven testing.
\item Mengintegrasikan pengujian ke GitHub Actions.
\end{itemize}
\textbf{Kata Kunci}: Cypress, TypeScript, E-Commerce, CI/CD}
\section{Minat}

\cvline{Teknologi}{Otomasi Pengujian, Open Source, Kecerdasan Buatan}
\cvline{Olahraga}{Lari, Bersepeda, Bulu Tangkis}
\section{Kegiatan Sukarela}

\cventry{Feb 2020–Sekarang}
        {Mentor Pengujian Perangkat Lunak}
        {Komunitas QA Indonesia}
        {\url{https://www.qaindonesia.id}}
        {}
        {\begin{itemize}
\item Menjadi mentor untuk QA junior dan mahasiswa melalui sesi belajar bulanan.
\item Membantu menyusun materi pengujian manual dan otomasi.
\item Berkontribusi pada forum diskusi dan webinar komunitas.
\end{itemize}}
    
\end{document}